% Source: https://artofproblemsolving.com/community/c3h1747745_a_phiery_nt_problem_from_me_to_you
% Stable identity: aops-topic-1747745-q01
\textbf{Which digits can occur.} The process starts at $n\ge10$ and stops at the first value below $10$, so the number put into $S$ is $\phi(m)$ for some $m\ge10$. For $m\ge3$, $\phi(m)$ is even: if $m$ has an odd prime factor $p$ with $p^a\parallel m$, then $\phi(m)$ is divisible by $\phi(p^a)=p^{a-1}(p-1)$, which is even; otherwise $m=2^a$ with $a\ge2$ and $\phi(m)=2^{a-1}$ is even. Also $\phi(m)\ge1$. So the digit is $2$, $4$, $6$ or $8$.

\textbf{The digit $2$ never occurs.} If $\phi(m)=2$ and a prime $p$ divides $m$, then $p-1$ divides $\phi(m)=2$, so $p\in\{2,3\}$. Writing $m=2^a3^b$: $\phi(m)=2$ forces $m\in\{3,4,6\}$ (for example $a\ge3$ gives $4\mid\phi(m)$, and $b\ge2$ gives $3\mid\phi(m)$). All of these are below $10$, so $\phi(m)=2$ never happens with $m\ge10$.

\textbf{The digits $4$, $6$, $8$ occur.} $\phi(10)=4$, $\phi(14)=6$, $\phi(15)=8$.

Hence $S=\{4,6,8\}$ and $|S|=\boxed{3}$.
